% Section 4.11, from lectures/L04/appendix/b_exercises.md.

\section{Exercises}
\label{c4:sec:exercises}

\begin{aside}
\textbf{How to check your work.} Every exercise below that asks you to write a VHDL module comes
with a self-checking testbench under \file{lectures/L04/exercises/}. Write your module in its
directory \file{exercises/<module>/}, with the entity name and the \textbf{port order} the exercise
states, and then run it with GHDL \textemdash{} see \secref{c2:sec:testbenches} for the three
commands.
\end{aside}

% ------------------------------------------------------------------------------------------
\exerciseset{Metastability concepts}

\exercise{Metastability in your own words}{Reasoning}
\label{c4:ex:meta}
Explain metastability in your own words. In your explanation, describe:
\begin{itemize}
  \item What happens inside a D flip-flop when its input changes too close to the active clock
    edge.
  \item Why the flip-flop's output can temporarily remain between a valid logic \code{0} and logic
    \code{1}.
  \item Why the output can take an unpredictable amount of time to settle.
  \item Why two downstream gates can interpret the unsettled output differently: one as \code{0},
    the other as \code{1}.
\end{itemize}

Explain why that disagreement can leave different parts of a digital circuit in inconsistent
states.

\exercise{Setup, hold, and why an asynchronous input cannot be avoided}{Reasoning}
\label{c4:ex:setup}
Setup time and hold time define a window around the active clock edge during which a flip-flop's
input must remain stable.

Explain:
\begin{itemize}
  \item Why a synchronously generated signal can normally be designed to meet those timing
    requirements.
  \item Why an asynchronous input, for example a mechanical push button, has no fixed relationship
    to the system clock.
  \item Why an asynchronous input can therefore change during the setup and hold window regardless
    of how the surrounding synchronous logic is designed.
\end{itemize}

Conclude why the metastability risk from an asynchronous input must be managed rather than simply
avoided.

% ------------------------------------------------------------------------------------------
\exerciseset{Designing a synchroniser}

\exercise{The chain, drawn and explained}{Circuit}
\label{c4:ex:chain}
Draw a two-flip-flop synchroniser for an asynchronous input. The circuit should contain an
asynchronous input, two D flip-flops connected in series, a shared system clock, and a synchronised
output taken from the second flip-flop.

Explain in your own words:
\begin{itemize}
  \item Why the first flip-flop is directly exposed to the asynchronous input.
  \item Why the first flip-flop can go metastable.
  \item How the interval between the first and second clock edges gives the first stage time to
    settle.
  \item Why the second flip-flop's output has a much lower probability of being metastable.
  \item Why it is the second stage, not the first, that the rest of the design should use.
\end{itemize}

A synchroniser does not make the probability of metastability exactly zero. Explain why it
nonetheless lowers the risk to an acceptably low level in most designs.

\note[Hint]Think about how much settling time the first flip-flop gets before its value is sampled
by the second flip-flop.

\exercise{A third stage ``just to be safe''}{Reasoning}
\label{c4:ex:third}
A colleague suggests using three synchroniser flip-flops on every asynchronous input in a slow
design clocked at a few megahertz, just to be safe.

A third stage buys another clock period for a metastable value to settle in, and costs another
clock cycle of latency on every synchronised input. At a few megahertz that extra clock period is
an enormous time in metastability terms, which is exactly why the suggestion sounds sensible, and
also why it may already be unnecessary.

\xpart{a} Reason through the trade-off in both directions. Under what circumstances does it come
out in favour of the third stage, and under what circumstances is the second stage already enough?
Your answer must hang on something other than the clock frequency, since frequency alone does not
settle it: say what it hangs on instead.

\xpart{b} \Secref{c4:sec:mtbf} arrived at MTBF from the settling time available per stage. Without
redoing that calculation, say which way each of these factors pushes the decision, and why:
\begin{itemize}
  \item How often the asynchronous input actually changes.
  \item The metastability characteristics of the FPGA family you are targeting.
  \item What a synchronisation failure would cost if it happened.
\end{itemize}

\xpart{c} Give your recommendation for an ordinary slow design of the kind found in this book. Then
describe a design, real or invented, where you would argue for the third stage instead, and say
which of the factors in \textbf{b)} does the work in that case.

\begin{aside}
\textbf{Check that your answer covers} the two things a third stage gives, another clock period of
settling time and a lower probability of an unsettled value reaching the functional logic, and the
two it costs, another cycle of latency and more flip-flops. If your answer to \textbf{a)} rests on
the clock frequency rather than on the required mean time between failures, read \secref{c4:sec:mtbf}
once more before you move on.
\end{aside}

\exercise{Assert asynchronously, release synchronously}{Reasoning}
\label{c4:ex:resetpattern}
Explain the reset pattern \textbf{assert asynchronously, release synchronously} from
\secref{c4:sec:resetsync}.

Describe the two halves separately:
\begin{itemize}
  \item Asynchronous assertion: the reset takes effect immediately, and does not wait for a clock
    edge.
  \item Synchronous release: the reset is removed only at an active clock edge, and all affected
    logic leaves reset in a controlled relationship with the clock.
\end{itemize}

Explain:
\begin{itemize}
  \item Why immediate assertion of reset can matter when the clock is stopped, when the clock is
    unstable, and when the system must enter its reset state without delay.
  \item Why releasing an asynchronous reset near a clock edge can violate a flip-flop's recovery or
    removal timing requirements.
  \item Why synchronising the release stops different flip-flops from leaving reset on different
    clock cycles.
\end{itemize}

% ------------------------------------------------------------------------------------------
\exerciseset{Building the circuit}

\exercise{\code{led_toggle_sync} by hand}{Circuit}
\label{c4:ex:circuit}
Build the \code{led_toggle_sync} circuit from \secref{c4:sec:worked} by hand in CircuitVerse.

\xpart{a} Add the external ports: the inputs \code{clock}, \code{reset_n} (asynchronous, active
low) and \code{button_n} (active low), plus the output \code{led}.

\xpart{b} Build the two-flip-flop reset synchroniser:
\begin{itemize}
  \item Connect the two flip-flops in series.
  \item Assert both stages asynchronously with the raw \code{reset_n} input.
  \item Let the reset propagate out of the chain synchronously.
  \item Use the synchronised reset output, \code{reset_s2_n}, throughout the rest of the circuit,
    and do not use the raw \code{reset_n} input directly outside the reset synchroniser.
\end{itemize}

\xpart{c} Build the three-flip-flop button synchroniser and the edge detector:
\begin{itemize}
  \item Use the first two flip-flops as a two-flip-flop synchroniser.
  \item Use the third flip-flop to store the previous synchronised button value.
  \item Compare the current and the previous synchronised value, and detect a button press when the
    current synchronised, active-low button value is \code{0} and the previous synchronised value
    is \code{1}.
  \item Use an AND gate with the current value inverted to generate \code{button_edge_s2}. The
    relationship is $\mathit{button\_edge\_s2} = \mathit{previous} \cdot \mathit{current}'$.
\end{itemize}

\xpart{d} Build the LED toggle logic: add a flip-flop storing \code{led_s}, reset it with
\code{reset_s2_n}, toggle it only when \code{button_edge_s2 = 1}, and connect the output \code{led}
to \code{led_s}.

\xpart{e} Simulate the whole design with a clock period of \code{1000 ms}. Assert and release
\code{reset_n}, press and release \code{button_n}, and observe the synchronised button stages,
\code{button_edge_s2} and \code{led_s}. Decide:
\begin{itemize}
  \item Whether the LED toggles at the first clock edge after the physical button press.
  \item How many clock cycles it takes for the button value to propagate through the synchroniser
    and the edge detector.
  \item Whether the observed latency agrees with the behaviour this chapter predicts.
\end{itemize}

% ------------------------------------------------------------------------------------------
\exerciseset{Debouncing and timing}

\exercise{What the chain does and does not do}{Reasoning}
\label{c4:ex:debounce}
You have now built and watched the circuit, so this is a question about what you saw.

The three-flip-flop chain in Exercise~\ref{c4:ex:circuit} does three jobs at once: it reduces the
risk of metastability, it stores the previous synchronised button value, and it detects a button
press. Explain why it does \textbf{not} by itself guarantee that the button is debounced.

\xpart{a} Separate the three things that are easy to confuse:
\begin{itemize}
  \item Synchronisation: turns an asynchronous input into a signal usable in the clock domain.
  \item Edge detection: generates a pulse when the sampled signal changes.
  \item Debouncing: stops several physical transitions from being read as several separate presses.
\end{itemize}
Which two of these does your circuit actually do?

\xpart{b} In your simulation one press gave a single clean toggle. That is a stronger claim than it
looks: it holds only under a particular condition about where the bounce landed relative to the
clock edges. State that condition exactly, in terms of what your circuit sampled and what it
therefore saw. Then explain why that condition is an accident of your clock period rather than a
property of your circuit. Be concrete about which of the two, the bounce or the sampling, your
design actually controls.

\xpart{c} Your simulation ran with a clock period of \code{1000 ms}; the FPGA runs at \code{50 MHz}.
A push button's contacts typically bounce on the order of 1~ms. Work out the consequences yourself:
\begin{itemize}
  \item How many times does a clock with a period of \code{1000 ms} sample the button during that
    millisecond of bounce?
  \item How many times does a \code{50 MHz} clock sample it?
  \item Every sampled transition looks like a new edge to your edge detector. From those two
    numbers, predict what the LED does on the board at a single physical press, and say why exactly
    the same circuit behaved impeccably in CircuitVerse.
\end{itemize}

\xpart{d} Describe one remedy of each kind: a hardware remedy (for example an RC low-pass filter
followed by a Schmitt-trigger input) and a purely digital remedy (for example requiring the
synchronised signal to hold a new level for some minimum number of clock cycles before accepting
it, or a counter that ignores further transitions for a fixed interval after a press
\textemdash{} which is what \chapref{c7:ch}'s timer would give you).

\xpart{e} Explain why the two-flip-flop synchroniser must be kept even once a debouncing method has
been added. Which problem does debouncing not solve?

% ------------------------------------------------------------------------------------------
\exerciseset{The synchronisers as modules of their own}

The worked example puts each synchroniser in a module of its own, so that every later design
instantiates the same two files instead of deriving them afresh. That is why both \chapref{c7:ch}'s
\code{walking_led} and \chapref{c8:ch}'s \code{fsm_led} ask you to copy these two in, and why
\chapref{c8:ch}'s \code{seq_detect_mealy}, which has no button, takes the reset synchroniser alone
and leaves the other one aside.

The two exercises that follow are you writing those modules. \Secref{c4:sec:resetsync} and
\ref{c4:sec:buttonsync} give you the logic; what they do not give you is the entity around it, and,
for the second, the step from a single button to a vector of them. Write each module yourself before
you open the worked example's copy, and then compare. Exercise~\ref{c4:ex:toggle2} then puts both
together into a working design.

\exercise{\code{reset_sync}}{VHDL}
\label{c4:ex:resetsync}
Write the reset synchroniser as a module named \code{reset_sync}.

The entity has:

\begin{booktable}{@{}p{6em}p{3.4em}p{5.2em}L@{}}
  \thead{Port} & \thead{Dir.} & \thead{Type} & \thead{Description} \\
  \midrule
  \code{clock} & in & \code{std_logic} & 50 MHz system clock. \\
  \code{reset_n} & in & \code{std_logic} & Active-low asynchronous reset, straight from the outside
    world. \\
  \code{reset_s2_n} & out & \code{std_logic} & Active-low synchronised reset, safe for the rest of
    the design to use. Goes low the moment \code{reset_n} does, with no clock edge involved, and
    returns high two rising edges after \code{reset_n} has. \\
\end{booktable}

\xpart{a} Implement \textbf{assert asynchronously, release synchronously}
(\secref{c4:sec:resetsync}):
\begin{itemize}
  \item An internal signal for the first stage, plus the output for the second.
  \item One process, sensitive to both \code{clock} and \code{reset_n}.
  \item At \code{reset_n = '0'}, drive both stages low, outside the clocked branch.
  \item At a rising clock edge, shift a constant \code{'1'} through the two stages.
\end{itemize}

\xpart{b} The value forced at reset and the value shifted in on the clock are each other's
opposites, \code{'0'} and \code{'1'}. Explain in one sentence why that is no contradiction, in
terms of what each branch is for.

\xpart{c} Verify the design with its testbench. It checks both halves of the pattern separately:
that a low \code{reset_n} pulls \code{reset_s2_n} low with no clock edge in between, and that
releasing \code{reset_n} leaves \code{reset_s2_n} low until two rising edges have passed.

\xpart{d} Answer in prose:
\begin{itemize}
  \item The output is called \code{reset_s2_n} rather than \code{reset_n_sync}. What does \code{s2}
    say, and what would have to change in this module for that name to become wrong?
  \item Every other module in this book takes \code{reset_s2_n} as its reset. This one takes the
    raw \code{reset_n}. Why must exactly one module in a design be the exception?
\end{itemize}

\note[Hint]The release is the interesting half. At reset the process does not care about the clock
at all, so both stages go low together; on the clock it does not care about \code{reset_n}, so the
\code{'1'} needs one edge per stage to reach the output. Two edges of latency at release is the
price of the whole pattern.

\modulefig[0.62]{lectures/L04/appendix/images/reset_sync.png}

\begin{selfcheck}
\textbf{Self-check.} Name your entity \code{reset_sync}, with the ports \code{clock},
\code{reset_n} (in) and \code{reset_s2_n} (out), declared in that order; its testbench is in
\file{exercises/reset_sync/}.
\end{selfcheck}

\exercise{\code{button_sync}}{VHDL}
\label{c4:ex:buttonsync}
Write the button synchroniser and edge detector as a module named \code{button_sync}.

Exercise~\ref{c4:ex:toggle2} handles two buttons by writing a two-bit synchroniser. This one is the
same circuit with the width left open: a generic decides how many buttons it serves, so that
\code{led_toggle_sync} can use it for one button and \code{fsm_led} for two without either file
having to change. See \secref{c4:sec:generics} for the syntax and for why this module has a generic
where \code{reset_sync} has none.

The entity has one generic:

\begin{booktable}{@{}p{4.6em}p{9.4em}p{4em}L@{}}
  \thead{Generic} & \thead{Type} & \thead{Default} & \thead{Description} \\
  \midrule
  \code{COUNT} & \code{natural range 1 to 3} & \code{1} & The number of buttons, and therefore the
    width of both vector ports. \\
\end{booktable}

and these ports:

\begin{booktable}{@{}p{7em}p{3.4em}p{11.4em}L@{}}
  \thead{Port} & \thead{Dir.} & \thead{Type} & \thead{Description} \\
  \midrule
  \code{clock} & in & \code{std_logic} & 50 MHz system clock. \\
  \code{reset_s2_n} & in & \code{std_logic} & Active-low, \textbf{already synchronised} reset. This
    module does not synchronise its own reset; the module in Exercise~\ref{c4:ex:resetsync} did
    that. \\
  \code{button_n} & in & \code{std_logic_vector(COUNT - 1 downto 0)} & Active-low asynchronous push
    buttons. \\
  \code{button_edge_s2} & out & \code{std_logic_vector(COUNT - 1 downto 0)} & A one-cycle pulse at
    every button's falling edge, two clock edges after the press. Every bit is independent of all
    the others. \\
\end{booktable}

\xpart{a} Build the three-stage chain from \secref{c4:sec:buttonsync}, now over vectors instead of
single bits:
\begin{itemize}
  \item Three internal \code{std_logic_vector(COUNT - 1 downto 0)} signals, one per stage.
  \item Stages 1 and 2 are the two-flip-flop synchroniser; stage 3 stores the previous value of
    stage 2.
  \item At reset, set all three stages to \code{(others => '1')}, not \code{(others => '0')}.
  \item Derive the output with a single concurrent assignment, and note that the same expression
    works unchanged on vectors as on bits.
\end{itemize}

\textbf{Keep this one in mind.} The first two flip-flops are a synchroniser in their entirety and
nothing else, and Exercise~\ref{c5:ex:sync} has you factor them out again as a module of their own,
generic in both width and the value they reset to. Writing them inline here first is what makes that
exercise a consolidation rather than a new idea.

\xpart{b} Explain why the reset value is all ones. All ones means ``released'', since the buttons
are active low, so a reset to all zeros starts the chain claiming that every button is already held
down. Follow what that costs:
\begin{itemize}
  \item With the buttons released and a reset to all zeros, does the module emit a false pulse as
    the released state propagates through the three stages? Walk through the output expression edge
    by edge before you answer, rather than guessing.
  \item Now the case that actually breaks. A user holds a button down at the moment the reset is
    released. What does the module report with a reset to all ones, and what does it report with all
    zeros? Which of the two is the behaviour you want, and why?
\end{itemize}

\xpart{c} Verify the design with its testbench. It instantiates your module \textbf{twice}, once
with \code{COUNT = 1} and once with \code{COUNT = 2}, and checks that a press gives exactly one
one-cycle pulse, that a held button gives no further pulses, that a release gives none at all, and
that pressing one button of a pair never pulses the other.

\xpart{d} Answer in prose:
\begin{itemize}
  \item \code{COUNT} is constrained to \code{1 to 3} rather than left as an unconstrained
    \code{natural}. What does that constraint buy, given that nothing in the architecture depends on
    the upper bound?
  \item This module is deliberately not merged with Exercise~\ref{c4:ex:resetsync}'s, even though a
    design that has a button always needs both. Name the design that justifies the split, and say
    what it would have had to do instead.
\end{itemize}

\note[Hint]Write the architecture for \code{COUNT = 1} in your head first, then check that every
line you wrote is already correct for a vector. \code{and}, \code{not} and the aggregate
\code{(others => '1')} all work element by element, so the generalisation should cost you no extra
code at all. If it does, that is worth another look.

\modulefig[0.86]{lectures/L04/appendix/images/button_sync.png}

\begin{selfcheck}
\textbf{Self-check.} Name your entity \code{button_sync}, with the generic \code{COUNT}
(\code{natural range 1 to 3}, default \code{1}), the inputs \code{clock}, \code{reset_s2_n},
\code{button_n} (\code{std_logic_vector(COUNT - 1 downto 0)}) and the output \code{button_edge_s2}
(\code{std_logic_vector(COUNT - 1 downto 0)}), declared in that order; its testbench is in
\file{exercises/button_sync/}.
\end{selfcheck}

% ------------------------------------------------------------------------------------------
\exerciseset{Putting them together into a design}

\exercise{\code{led_toggle_sync2}}{VHDL}
\label{c4:ex:toggle2}
Now make \chapref{c3:ch}'s \code{led_toggle} metastability-safe, in VHDL, as a new module named
\code{led_toggle_sync2}.

\chapref{c3:ch}'s \code{led_toggle} toggled two LEDs from two push buttons, but fed their raw,
asynchronous signals straight into the edge detector \textemdash{} logically correct, but unsafe on
real hardware. Here you rebuild it with every asynchronous input protected by this chapter's
two-flip-flop synchroniser.

The entity has:

\begin{booktable}{@{}p{5.4em}p{3.4em}p{11.4em}L@{}}
  \thead{Port} & \thead{Dir.} & \thead{Type} & \thead{Description} \\
  \midrule
  \code{clock} & in & \code{std_logic} & 50 MHz system clock. \\
  \code{reset_n} & in & \code{std_logic} & Active-low asynchronous reset. \\
  \code{button_n} & in & \code{std_logic_vector(1 downto 0)} & Two active-low asynchronous push
    buttons. \\
  \code{led} & out & \code{std_logic_vector(1 downto 0)} & Two LEDs; \code{led(i)} toggles once per
    press of \code{button_n(i)}. \\
\end{booktable}

\xpart{a} Synchronise every asynchronous input:
\begin{itemize}
  \item Apply the two-flip-flop reset synchroniser from \secref{c4:sec:resetsync}.
  \item Give the buttons a two-flip-flop synchroniser plus a third stage for edge detection, as in
    \secref{c4:sec:buttonsync} \textemdash{} now over a \code{std_logic_vector(1 downto 0)}, so that
    both buttons are synchronised in parallel.
  \item Never let the raw \code{reset_n} or \code{button_n} reach the toggle logic; use only the
    synchronised reset and the per-button edge pulses.
\end{itemize}

\note[Hint]You can put the synchroniser inline, or factor it out into a subcomponent taking a
two-bit \code{button_n}, with the instantiation syntax from \secref{c2:sec:submodules}. Inline is
the shorter route, and then the build is a single file:

\begin{shell}
ghdl -a --std=93 led_toggle_sync2.vhd led_toggle_sync2_tb.vhd
\end{shell}

If you factor it out instead, copy in the \file{reset_sync.vhd} and \file{button_sync.vhd} you wrote
for Exercise~\ref{c4:ex:resetsync} and \ref{c4:ex:buttonsync} (and instantiate \code{button_sync}
with \code{COUNT = 2}), and name them first on the \code{ghdl -a} line, before your own module, so
that GHDL sees them before the file that instantiates them:

\begin{shell}
ghdl -a --std=93 reset_sync.vhd button_sync.vhd \
                 led_toggle_sync2.vhd led_toggle_sync2_tb.vhd
\end{shell}

See \secref{c2:sub:multifile}. Neither module is handed out anywhere in the course repository,
because writing them is Exercise~\ref{c4:ex:resetsync} and \ref{c4:ex:buttonsync}.

\xpart{b} Build the toggle logic:
\begin{itemize}
  \item One LED-state flip-flop per LED.
  \item At the synchronised reset, turn both LEDs off.
  \item Toggle \code{led(i)} only when button \code{i}'s synchronised falling-edge pulse fires.
  \item Confirm that \code{led(i)} responds to \code{button_n(i)} alone, never to the other button.
\end{itemize}

\xpart{c} Verify the design with its testbench. The testbench presses each button in turn and checks
that reset turns both LEDs off, that no press means no toggle, that pressing button 0 toggles
\textbf{only} \code{led(0)} once the synchroniser latency has elapsed, that holding that button down
does not toggle it again, and that pressing button 1 then toggles only \code{led(1)} and leaves
\code{led(0)} as it was.

Pay attention to the \textbf{latency}: the testbench checks the \emph{exact} edge, not merely that
the LED eventually settles correctly. A press has to reach the toggle logic through two synchroniser
flip-flops, so \code{led} is still unchanged after the first and the second rising edge and toggles
at the third. Count the flip-flops a press travels through and convince yourself that the third edge
is the right one, and that a defensive extra pipeline stage would now be caught rather than
tolerated.

\xpart{d} Answer in prose, comparing this design with your \code{led_toggle} from \chapref{c3:ch}:
\begin{itemize}
  \item What does the synchroniser change about pressing a real, bouncing button?
  \item Why was the raw version unsafe to connect to a physical button at \code{50 MHz}?
  \item Your testbench drives clean, perfectly timed transitions. Which of this chapter's two
    problems, metastability or contact bounce, can a testbench therefore never demonstrate to you?
\end{itemize}

\modulefig[0.68]{lectures/L04/appendix/images/led_toggle_sync2.png}

\begin{selfcheck}
\textbf{Self-check.} Name your entity \code{led_toggle_sync2}, with the inputs \code{clock},
\code{reset_n}, \code{button_n} (\code{std_logic_vector(1 downto 0)}) and the output \code{led}
(\code{std_logic_vector(1 downto 0)}), declared in that order; its testbench is in
\file{exercises/led_toggle_sync2/}. It presses each button in turn and checks that \code{led(i)}
toggles on \code{button_n(i)}'s falling edge, that it does so at the right clock edge rather than
merely settling on the right value eventually, and that it never fires the wrong LED.
\end{selfcheck}
